\section{СИМУЛЯЦИЯ ПОВЕДЕНИЯ АВТОНОМНЫХ АГЕНТОВ В ГОРОДСКОЙ СРЕДЕ}

Сравнительная оценка методов глобального планирования требует воспроизводимой среды, в которой можно 
одинаковым образом запускать различные алгоритмы на одних и тех же входных данных и измерять выбранные 
показатели качества. Натурный эксперимент с реальными доставочными агентами для этой цели плохо 
приспособлен: он дорог, плохо повторяем и не позволяет в разумные сроки перебрать достаточное количество
конфигураций дорожной сети и плотности агентов. Поэтому в настоящей работе предложен и реализован 
симулятор, выступающий основным инструментом экспериментального исследования.

Симулятор должен решать следующие задачи:

\begin{itemize}
\item воспроизводить движение группы агентов по взвешенному графу
      дорожной сети с узкими местами;
\item предоставлять единый интерфейс для подключения как методов
      планирования, заранее учитывающих возможные конфликты, так и
      методов непосредственного разрешения конфликтов в узких местах;
\item обеспечивать загрузку графов, построенных по реальным фрагментам
      городской среды, и генерацию конфигураций агентов с заданными
      параметрами;
\item вычислять метрики качества и сохранять
      результаты прогонов в виде, пригодном для последующего
      статистического сравнения.
\end{itemize}

В настоящего разделе будет описано устройство разработанного симулятора на уровне, достаточном для 
понимания того, как именно получены экспериментальные результаты, приведённые далее.

\subsection{Архитектура симулятора}

\subsection{Компоненты симулятора}

\subsection{Получение исходных данные}

\subsection{Используемые технологии}